\documentclass[11pt]{article}


\setlength{\parindent}{0pt}
\usepackage{xltxtra}
\usepackage{hyperref}
\hypersetup{hidelinks}
\usepackage{url}
\urlstyle{tt}
\usepackage{xcolor}
\definecolor{CVBlue}{RGB}{23,110,191}
\usepackage{calc}
\usepackage{graphicx}
\usepackage{tikz}
\usetikzlibrary{calc}
\usepackage{fontspec}
\usepackage{xeCJK}
\usepackage{enumitem}
\CJKsetecglue{} %% 取消中文与数字之间的间隙


%% 主文档字体设置
\setmainfont[
    Path = fonts/Main/,
    Extension = .otf,
    BoldFont = texgyretermes-bold.otf, % 加粗字体
]{texgyretermes-regular.otf} % 正文字体

% 中文字体设置
\setCJKmainfont[
    Path = fonts/hansans/,
    Extension = .ttf,
    BoldFont = NotoSansSC-Bold.ttf, % 加粗字体
]{NotoSansSC-Regular.otf} % 正文字体


\usepackage{relsize}
\usepackage{xspace}

% 使用 fontawesome（部分图标）
\usepackage{fontawesome} 

% A4纸，上下左右边距
\usepackage[
    a4paper,
    left=1.2cm,
    right=1.2cm,
    top=1.5cm,
    bottom=1cm,
    nohead
]{geometry}

\renewcommand{\baselinestretch}{1.5} % 行间距设为1.5

\usepackage{titlesec}
\usepackage{enumitem}
\setlist{noitemsep} % 取消列表项间的额外间距
%\setlist{nosep} % 取消所有垂直间距
\setlist[itemize]{topsep=0.25em, leftmargin=*}
\setlist[enumerate]{topsep=0.25em, leftmargin=*}

% --- 用于控制【不同项目之间】的垂直距离 ---
\newlength{\interProjectSpacing}
\setlength{\interProjectSpacing}{0.9em} % <--- 在此调整项目之间的距离
\newcommand{\projectsep}{\vspace{\interProjectSpacing}}

% --- 用于控制【项目标题】与下方【项目描述】的距离 ---
\newlength{\intraProjectTitleSep}
\setlength{\intraProjectTitleSep}{0.4em} % <--- 在此调整标题和描述的距离
\newcommand{\titlebreak}{\\[\intraProjectTitleSep]}

% --- 用于控制【项目描述】与下方【要点列表】的距离 ---
\newlength{\intraProjectListTopSep}
\setlength{\intraProjectListTopSep}{0.2em} % <--- 在此调整描述和列表的距离

% =======================================================================


\titleformat{\section}         % 定制 \section 命令 
{\large\bfseries\raggedright} % 将 section 标题设置为大号、粗体且左对齐
{}{0em}                      % 可用于为所有 section 添加前缀（如“章节...”）
{}                           % 可用于在标题前插入代码
[{\color{CVBlue}\titlerule}]  % 在标题后插入一条横线
\titlespacing*{\section}{0cm}{*1.6}{*1.2}



\begin{document}
\pagenumbering{gobble}

%%%% 利用tikz来定位照片
\begin{tikzpicture}[remember picture, overlay] 
    \node[anchor = north east] at ($(current page.north east)+(-2cm,-1.2cm)$) {\includegraphics[height=3cm]{avatar.jpg}};
  \end{tikzpicture}%
  %%%% 利用tikz来定位学校Logo，这里只在第一页显示
  \begin{tikzpicture}[remember picture, overlay] 
    \node[anchor = north west] at ($(current page.north west)+(0.5cm,+1.0cm)$) {\includegraphics[height=6cm]{zju.png}};
  \end{tikzpicture}%
\centerline{\LARGE\bfseries{郑煜}} 

\centerline{\normalsize{\faPhone\ 173-5778-2028 \quad \faEnvelopeO\ \href{mailto:790245083@qq.com}{790245083@qq.com}}} 


    
\section{\makebox[\widthof{\faGraduationCap}][c]{\color{CVBlue}\faGraduationCap}\ 教育背景}    
\textbf{浙江大学} \hfill 2024.9 -- 2027.6\\[0.5em] % 标题和正文间加一点距离
计算机科学与技术 硕士\quad  
\begin{itemize}[nosep]
    \item 相关课程：《cuda编程》、《机器学习》、《数据挖掘》
\end{itemize}

\section{\makebox[\widthof{\faUsers}][c]{\color{CVBlue}\faUsers}\ 项目经历}

% --- 第一个项目 ---
% 将标题行末尾的 \\ 替换为 \titlebreak 命令
\textbf{自主实现类 ChatGPT 大模型训练} \hfill 2025.08 -- 至今 \titlebreak
项目描述：从零构建端到端大语言模型训练系统，完整覆盖类 ChatGPT 的核心技术栈，实现从数据处理到模型训练的全流程自主开发与优化。
% 在 itemize 的选项中，使用 topsep=\intraProjectListTopSep 来控制上边距
\begin{itemize}[nosep, topsep=\intraProjectListTopSep]
    \item \textbf{训练}：基于 PyTorch 构建端到端大模型训练流水线，完成 FineWeb 数据清洗、去重与分片处理，覆盖 Tokenizer 构建、预训练（Pre-training）、监督微调（SFT）及基于 PPO 的 RLHF 对齐训练全流程。
在 8×A100 集群环境下完成 1.05M tokens 规模训练，采用混合精度与系统化超参数调优优化收敛稳定性与算力利用率。
    \item \textbf{核心产出}：最终在 880M 参数规模下实现 Core Score 相较 GPT-2 提升 2%。
\end{itemize}

% 使用 \projectsep 命令来分隔两个项目
\projectsep

% --- 第二个项目 ---
\textbf{开源大模型技术知识(RAG)系统} \hfill 2025.02 -- 至今 \titlebreak
项目描述：基于 Qwen 系列大模型,面向 GitHub 上 LLM / GPT / Transformer 相关开源项目构建检索增强生成（RAG）问答系统.
\begin{itemize}[nosep, topsep=\intraProjectListTopSep]
    \item \textbf{文档构建}：基于 GitHub API 批量采集 Star >100 的 LLM 相关开源项目 README 文档，构建领域技术语料；设计规则清洗管线完成格式标准化与噪声过滤，并引入大模型进行语义级重构与冗余压缩，最终完成结构化分片以适配向量检索与 RAG 系统。
    \item \textbf{RAG召回}：文档向量化采用 Qwen Embedding v4，向量数据存储于 Milvus 数据库中。RAG 调用策略采用 规则判断与 LLM 自判断的混合机制，根据问题类型和模型置信度动态决定是否触发检索增强生成。
    \item \textbf{可视化GUI}：前端使用 \textbf{React} 框架构建了\textbf{后台管理的可视化面板 (Admin Panel)}。可一键下达挖坑指令、动态调整Agent的\textbf{指令}、并实时监控数据。
\end{itemize}

\section{\makebox[\widthof{\faCogs}][c]{\color{CVBlue}\faCogs}\ 技术栈}
\begin{itemize}[nosep]
    \item \textbf{编程语言：} \textbf{Go}, Python, C++
    \item \textbf{开发工具：} SSH, Git, Vim, MakeFile,LaTex
    \item \textbf{操作系统：} Linux
\end{itemize}
\section{\makebox[\widthof{\faGraduationCap}][c]{\color{CVBlue}\faList}\ 获奖情况}
\begin{itemize}
    \item CC98年度用户 \hfill 2023.12
    
\end{itemize}
    
\section{\makebox[\widthof{\faInfo}][c]{\color{CVBlue}\faInfo}\ 其他}
\begin{itemize}[parsep=0.5ex]
    \item \textbf{技术博客：} \href{https://maksymilan.github.io/}{https://maksymilan.github.io/}
    \item \textbf{GitHub：} \href{https://github.com/maksymilan}{https://github.com/maksymilan} 
    \item \textbf{英语水平：} CET-4, CET-6
\end{itemize}
\end{document}
